\begin{sidewaystable}
\footnotesize\setlength{\tabcolsep}{4pt}
\caption{INT8 accuracy under different range-setting procedures}\label{tab:calibmethods}
\begin{tabular}{@{}lrrrrrr@{}}
\toprule
Model & FP32 & \multicolumn{3}{c}{INT8, toolchain entropy calibration} & INT8, order-independent & FP8 \\
\cmidrule{3-5}
 & & full set & subsets (range) & max $|\Delta|$ & entropy calibration & \\
\midrule
Faster R-CNN R50-FPN & 36.95 & 36.71 & 36.70--36.72 & 0.01 & 36.74 & 36.58 \\
YOLOv10-S & 46.03 & 43.79 & 43.82--45.31 & 1.52 & 35.80 & 45.62 \\
YOLOv10-X & 54.00 & 47.79 & 47.79--50.45 & 2.66 & 37.07 & 53.85 \\
EfficientNet-B0 & 65.79 & 24.82 & 40.09--41.70 & 16.88 & 0.09 & 65.40 \\
DenseNet-121 & 62.03 & 57.36 & 58.31--60.44 & 3.09 & 60.01 & 61.71 \\
EfficientNet-B0, convolution inputs only & 65.79 & 23.84 & -- & -- & 54.91 & -- \\
\botrule
\end{tabular}
\footnotetext{\textbf{What to look for:} INT8 accuracy depends strongly on how the ranges are set: the order-independent procedure is much worse for the YOLOv10 networks and EfficientNet-B0 but better for DenseNet-121. FP8, whose range is simply the maximum, stays within half a point of FP32 for all five networks. Toolchain: incremental entropy calibration of ModelOpt on ONNX Runtime (one image per step), full calibration set in stored order (main results) and seeded random subsets, which differ in their first image (Table~\ref{tab:calibvar}). Order-independent: global 2048-bin histogram with NVIDIA's reference KL search (Supplementary Section~\ref{sec:supp-quant}), applied to the same quantizers. Convolution inputs only: quantizers on the inputs of convolutions and matrix multiplications, as in~\cite{wu2020integer}. COCO val2017 box mAP and ImageNetV2 top-1 (\%); FP8: max calibration.}
\end{sidewaystable}
